\documentclass{beamer}

% --- Theme Selection ---
\usetheme{Madrid}
\usecolortheme{whale}

% --- Packages ---
\usepackage[utf8]{inputenc}
\usepackage{graphicx}
\usepackage{amsmath}
\usepackage{listings}
\usepackage{hyperref}

% --- Code Snippet Styling ---
\lstset{
    basicstyle=\ttfamily\scriptsize,
    breaklines=true,
    backgroundcolor=\color{gray!10},
    frame=single,
    keywordstyle=\color{blue},
    commentstyle=\color{green!50!black}
}

% --- Metadata ---
\title[Visual Docking Station]{Autonomous Visual Docking Station}
\subtitle{Precision Navigation with ROS 2 Jazzy \& Gazebo Harmonic}
\author{Presented by Project Team}
\date{March 2026}

\begin{document}

% --- Title Slide ---
\begin{frame}
    \titlepage
\end{frame}

% --- Slide 1: Objectives ---
\begin{frame}{Project Objectives}
    \begin{itemize}
        \item \textbf{Autonomous Homing:} Robot must locate and navigate to a docking station from any room position.
        \item \textbf{Visual Perception:} Utilize ArUco markers for robust 6-DOF pose estimation.
        \item \textbf{Precision Control:} Implement a state-machine driven controller for smooth approach and docking.
        \item \textbf{Simulation Fidelity:} High-accuracy physics modeling in Gazebo Harmonic.
    \end{itemize}
\end{frame}

% --- Slide 2: System Architecture ---
\begin{frame}{System Architecture}
    \begin{columns}
        \begin{column}{0.5\textwidth}
            \textbf{Core Stack:}
            \begin{itemize}
                \item OS: Ubuntu 24.04 (Noble)
                \item ROS 2 Jazzy Jalisco
                \item Gazebo Harmonic
            \end{itemize}
        \end{column}
        \begin{column}{0.5\textwidth}
            \textbf{Key ROS Nodes:}
            \begin{itemize}
                \item \texttt{aruco\_detector\_node}: OpenCV-based vision
                \item \texttt{dock\_controller}: 3-phase controller
                \item \texttt{ros\_gz\_bridge}: Middleware bridge
            \end{itemize}
        \end{column}
    \end{columns}
\end{frame}

% --- Slide 3: Environment Design ---
\begin{frame}{Simulation Environment}
    \begin{itemize}
        \item \textbf{Room Layout:} 8m $\times$ 6m interior with alcove.
        \item \textbf{Docking Target:} ArUco Marker ID 0 (DICT\_4X4\_50).
        \item \textbf{Lighting:} Point lights optimized for camera contrast.
    \end{itemize}
    \begin{block}{The Alcove}
        The docking platform is recessed 1.9m into an alcove to simulate realistic indoor constraints.
    \end{block}
\end{frame}

% --- Slide 4: Robot Modeling (SDF) ---
\begin{frame}{Robot Modeling \& Physics}
    \begin{itemize}
        \item \textbf{Kinematics:} Differential-drive system.
        \item \textbf{Hardware:}
            \begin{itemize}
                \item Dual motorized wheels.
                \item Low-friction ball caster ($\mu = 0$).
                \item Forward RGB Camera (640$\times$480).
            \end{itemize}
        \item \textbf{Logic:} Defined in native SDF version 1.9 for direct Gazebo compatibility.
    \end{itemize}
\end{frame}

% --- Slide 5: Perception Logic ---
\begin{frame}[fragile]{Perception: ArUco Tracking}
\begin{lstlisting}[language=Python]
# OpenCV Detection Flow
detector = aruco.ArucoDetector(dictionary, parameters)
corners, ids, _ = detector.detectMarkers(image)
if ids is not None:
    _, rvec, tvec = cv2.solvePnP(obj_pts, corners[0], K, D)
    # Publish Pose: depth=tvec[2], lateral=tvec[0]
\end{lstlisting}
\begin{itemize}
    \item Modern OpenCV 4.8+ non-deprecated API.
    \item PnP algorithm provides distance and angle in real-time.
\end{itemize}
\end{frame}

% --- Slide 6: Control Logic ---
\begin{frame}{State Machine Controller}
    \textbf{1. SEARCHING:}
    \begin{itemize}
        \item Rotate slowly until marker is identified.
    \end{itemize}
    \textbf{2. APPROACHING (Visual Servoing):}
    \begin{itemize}
        \item Angular: $K_a \times \text{lateral\_error}$
        \item Linear: $K_l \times (\text{distance} - \text{stop\_dist})$
    \end{itemize}
    \textbf{3. DOCKED:}
    \begin{itemize}
        \item Stop when distance $< 0.55\text{m}$.
    \end{itemize}
\end{frame}

% --- Slide 7: Challenges Overcome ---
\begin{frame}{Critical Implementation Challenges}
    \begin{alertblock}{Physics Stall}
        Fixed caster wheel friction was originally too high, preventing motion. Fixed by setting $\mu=0$ in SDF.
    \end{alertblock}
    \begin{block}{Topic Scoping}
        Gazebo Harmonic scopes topics under \texttt{/model/}. Configured explicit bridge remapping to allow ROS compatibility.
    \end{block}
\end{frame}

% --- Slide 8: Conclusion ---
\begin{frame}{Conclusion \& Performance}
    \begin{itemize}
        \item \textbf{Speed:} Rapid detection using DICT\_4X4\_50.
        \item \textbf{Precision:} High-reliability docking accuracy.
        \item \textbf{Scalability:} Clean, modular ROS 2 node architecture.
    \end{itemize}
    \begin{center}
        \textbf{Thank You!} \\
        Questions?
    \end{center}
\end{frame}

\end{document}
